\scanpage{121}
près et les structures de polarisation sur $G^{\unsure{\circ|0}}$ se \uline{déterminent mutuellement}.

\deleted{\unsure{Donc}} [À voir: l'ens des autom de $(G^{\unsure{\circ|0}}_{\mathbb R},\varepsilon_{\mathbb R},j_{\mathbb R})$ voire $(G^{\unsure{\circ|0}}_{\unsure{Q}},\varepsilon,\unsure{j|\mu})$, voire $(G,\varepsilon,\unsure{j|\mu})$ opère-t-il transitivement sur $\operatorname{Pol}(G)$ ? \deleted{Par exemple} est-il possible \deleted{\unsure{ici}} de distinguer \deleted{\unsure{au moyen de}} \inserted{une} \deleted{\unsure{toujours, y a-t-il}} polarisation ou un paquet de polarisations préférentielles ?]

Prenons p.ex. une catégorie de Motifs géométriques sur $\mathbb F_q$, $q=p^n$, \deleted{\unsure{et}} Alors \deleted{\unsure{donc}} il y a une classe de conjugaison d'homs
\diagram{\begin{tikzcd}
\unsure{i}:G_{m\mathbb C}\arrow[r] & G_{\mathbb C}
\end{tikzcd}}
distinguée, avec
\[
\bar i i=j,\qquad\varepsilon i=\mathrm{id}
\]
\margin{N.B. $i$ et \unsure{les conjugués} \inserted{\unsure{distincts}} sur $Q$, (\unsure{au nb} \deleted{\unsure{de}} $\unsure{2n}$, si $[K:Q]=n$ \unsure{engendrent} $\operatorname{Hom}(G_{m\mathbb C},G^{\unsure{\circ|0}}_{\mathbb C})$}
définie, \deleted{par} si $G$ commutatif: $M\subset L$ \margin{$L$ quadratique tot. imag., $K$ tot. réelle,} par la
\figure{\begin{tikzpicture}[every node/.style={align=left}]
\node[draw,ellipse,inner sep=5pt] (f) at (5,1.7) {forme};
\node[draw,rounded corners=12pt,text width=4cm,inner sep=5pt] (a) at (0,0) {\deleted{\unsure{est}}\par \uuline{\unsure{caveat}} $\longrightarrow$ cette forme \deleted{\unsure{en général}} n'est pas tjrs ! à valeurs entières \unsure{même si}};
\node[draw,rounded corners=15pt,text width=4cm,inner sep=5pt] (b) at (0,-3.7) {$\unsure{M}$ si $j=0$, ex: prendre $q=p\ne2$\par $L=Q(\unsure{\sqrt{-p}|\sqrt p})$, tel que nb de classes de $L$ soit $1$, p.ex.\par $p:3,7,11,\ldots$};
\draw (a.north east)--(f.west);
\draw (a.south)--(b.north);
\end{tikzpicture}}
\[
\lambda\longmapsto\frac1n v_{\mathfrak p}(\lambda),
\]
si $\mathfrak p$ est un idéal premier de $L$ sur $p$.
\[
[\ v_p(\lambda)v_p(\bar\lambda)=v_p(\lambda\bar\lambda)=v_p(q^{j^*(\lambda)})=j^*(\lambda);\quad v_p(q)=1\ ].
\]
En fait, il convient de ``normaliser'' $G_{\mathbb C}$ \inserted{\fbox{sur $Q$ !}} \unsure{choix} d'un tore; (NB. Un autom de $G^{\unsure{\circ|0}}$ qui \deleted{induit} respecte $i$ est l'identité !). La structure d'effectivité sur $M$ est déterminée grâce à $i$, celle des $M_k$ ($k$ ext. de $Q$) en particulier de $M_{\mathbb R}$, $M_{\mathbb C}$, par la considération des systèmes de leurs conjugués
\mspage{120}
